\section{Conexión con DBeaver}

\subsection{Inicio de DBeaver}

DBeaver Community 26.2.0 se instaló como aplicación de escritorio. En el
primer inicio se presentó un cierre inesperado con un error nativo
\texttt{SIGBUS} en \texttt{libgtk-3.so.0}; el registro indicó que se estaba
ejecutando Java 25.0.3. Para continuar, se instaló OpenJDK 21 y DBeaver se
inició usando ese entorno, con GTK mediante X11.

\begin{lstlisting}[style=terminal]
$ sudo apt install -y openjdk-21-jre
$ GDK_BACKEND=x11 dbeaver -clean \\
    -vm /usr/lib/jvm/java-21-openjdk-amd64/bin/java
\end{lstlisting}

\subsection{Parámetros de conexión}

Se creó la conexión \texttt{Practica01} con los siguientes parámetros:

\begin{table}[H]
  \centering
  \rowcolors{2}{unamblue!6}{white}
  \begin{tabular}{ll}
    \toprule
    \textbf{Parámetro} & \textbf{Valor} \\
    \midrule
    Nombre de la conexión & Practica01 \\
    Servidor & \texttt{127.0.0.1} \\
    Puerto publicado & \texttt{5433} \\
    Base de datos & \texttt{postgres} \\
    Usuario & \texttt{postgres} \\
    Contraseña & Credencial temporal oculta \\
    \bottomrule
  \end{tabular}
  \caption{Parámetros usados para conectar DBeaver con el contenedor de la práctica.}
\end{table}

La validación se realizó desde un editor SQL de DBeaver mediante la siguiente
consulta:

\begin{lstlisting}[style=terminal]
SELECT
    current_database(),
    current_user,
    version(),
    inet_server_addr(),
    inet_server_port();
\end{lstlisting}

El resultado mostró la base \texttt{postgres}, el usuario \texttt{postgres} y
PostgreSQL 18.6. La interfaz mostró \texttt{localhost:5433}; el valor 5432 de
\texttt{inet\_server\_port()} corresponde al puerto interno del contenedor.

\evidencia{04_dbeaver_conectado.png}{Conexión activa de DBeaver mediante \texttt{localhost:5433} y resultado de la consulta de validación.}
